\documentclass[11pt,a4paper]{article}

\usepackage[utf8]{inputenc}
\usepackage[T1]{fontenc}
\usepackage[french]{babel}
\usepackage[a4paper,margin=2.3cm]{geometry}
\usepackage{amsmath,amssymb,amsthm,mathtools}
\usepackage{lmodern}
\usepackage{enumitem}
\usepackage{hyperref}
\usepackage{xcolor}
\usepackage{fancyhdr}
\usepackage{stmaryrd}
\usepackage{mathrsfs}

\hypersetup{
    colorlinks=true,
    linkcolor=black,
    urlcolor=blue,
    citecolor=black
}

\setlength{\parindent}{0pt}
\setlength{\parskip}{0.6em}

\pagestyle{fancy}
\fancyhf{}
\fancyhead[L]{Algèbre linéaire --- Chapitre 4}
\fancyhead[R]{Familles libres, génératrices et bases}
\fancyfoot[C]{\thepage}

% Environnements
\newtheorem{definition}{Définition}[section]
\newtheorem{theoreme}[definition]{Théorème}
\newtheorem{proposition}[definition]{Proposition}
\newtheorem{corollaire}[definition]{Corollaire}
\newtheorem{lemme}[definition]{Lemme}
\theoremstyle{remark}
\newtheorem{remarque}[definition]{Remarque}
\newtheorem{exemple}[definition]{Exemple}
\newtheorem*{methode}{Méthode}
\newtheorem*{idee}{Idée}

% Commandes
\newcommand{\R}{\mathbb{R}}
\newcommand{\N}{\mathbb{N}}
\newcommand{\Z}{\mathbb{Z}}
\newcommand{\Q}{\mathbb{Q}}
\newcommand{\C}{\mathbb{C}}
\newcommand{\K}{\mathbb{K}}
\newcommand{\Vect}{\mathrm{Vect}}

\begin{document}

\begin{center}
    {\LARGE \textbf{Chapitre 4 --- Familles libres, génératrices et bases}}\\[0.5em]
    {\large Décrire un espace vectoriel à l’aide d’une famille de vecteurs}
\end{center}

\hrule
\vspace{1em}

\tableofcontents

\newpage

\section{Pourquoi ce chapitre est-il central ?}

Dans le chapitre précédent, on a vu qu’un espace vectoriel est un ensemble dans lequel on peut former des combinaisons linéaires. La question naturelle devient alors :

\begin{center}
\textit{comment décrire efficacement tous les vecteurs d’un espace vectoriel ?}
\end{center}

Pour cela, trois notions sont fondamentales :
\begin{itemize}
    \item les \textbf{familles génératrices}, qui permettent de fabriquer tous les vecteurs ;
    \item les \textbf{familles libres}, qui garantissent qu’il n’y a pas de redondance ;
    \item les \textbf{bases}, qui réalisent le meilleur compromis possible entre ces deux idées.
\end{itemize}

\begin{remarque}
Ce chapitre est au cœur de l’algèbre linéaire : presque toute la suite du cours repose sur la capacité à reconnaître une famille libre, une famille génératrice ou une base.
\end{remarque}

\section{Rappels sur les combinaisons linéaires}

\begin{definition}
Soient \(u_1,\dots,u_p\) des vecteurs d’un espace vectoriel \(E\).  
Une \textbf{combinaison linéaire} de \(u_1,\dots,u_p\) est tout vecteur de la forme
\[
\lambda_1u_1+\cdots+\lambda_pu_p,
\qquad
\text{avec } \lambda_1,\dots,\lambda_p\in\K.
\]
\end{definition}

\begin{definition}
Le sous-espace engendré par \(u_1,\dots,u_p\), noté
\[
\Vect(u_1,\dots,u_p),
\]
est l’ensemble de toutes les combinaisons linéaires de ces vecteurs :
\[
\Vect(u_1,\dots,u_p)=
\left\{
\lambda_1u_1+\cdots+\lambda_pu_p \mid \lambda_1,\dots,\lambda_p\in\K
\right\}.
\]
\end{definition}

\begin{exemple}
Dans \(\R^2\),
\[
\Vect((1,0),(0,1))=\R^2.
\]
Dans \(\R^3\),
\[
\Vect((1,0,0),(0,1,0))=\{(x,y,0)\mid x,y\in\R\}.
\]
\end{exemple}

\section{Familles génératrices}

\subsection{Définition}

\begin{definition}
Soit \(E\) un espace vectoriel. Une famille \((u_i)_{i\in I}\) de vecteurs de \(E\) est dite \textbf{génératrice} de \(E\) si tout vecteur de \(E\) est combinaison linéaire de vecteurs de cette famille.

Autrement dit, si la famille est finie \((u_1,\dots,u_p)\), elle est génératrice de \(E\) si
\[
E=\Vect(u_1,\dots,u_p).
\]
\end{definition}

\begin{idee}
Une famille génératrice est une famille à partir de laquelle on peut \textbf{fabriquer tout l’espace}.
\end{idee}

\begin{exemple}
Dans \(\R^2\), la famille
\[
\big((1,0),(0,1)\big)
\]
est génératrice, car tout \((x,y)\in\R^2\) s’écrit
\[
(x,y)=x(1,0)+y(0,1).
\]
\end{exemple}

\begin{exemple}
Dans \(\R^2\), la famille
\[
\big((1,0),(0,1),(1,1)\big)
\]
est aussi génératrice. Mais elle contient un vecteur redondant, car
\[
(1,1)=(1,0)+(0,1).
\]
\end{exemple}

\subsection{Comment montrer qu’une famille est génératrice ?}

\begin{methode}
Pour montrer qu’une famille \((u_1,\dots,u_p)\) est génératrice d’un espace \(E\), on prend un vecteur arbitraire \(x\in E\) et on montre qu’il existe des scalaires \(\lambda_1,\dots,\lambda_p\) tels que
\[
x=\lambda_1u_1+\cdots+\lambda_pu_p.
\]
\end{methode}

\begin{exemple}
Montrons que \(\big((1,1),(1,-1)\big)\) est génératrice de \(\R^2\).

Soit \((x,y)\in\R^2\). Cherchons \(\lambda,\mu\in\R\) tels que
\[
(x,y)=\lambda(1,1)+\mu(1,-1).
\]
Cela donne le système
\[
\begin{cases}
\lambda+\mu=x,\\
\lambda-\mu=y.
\end{cases}
\]
En additionnant et en soustrayant,
\[
\lambda=\frac{x+y}{2},
\qquad
\mu=\frac{x-y}{2}.
\]
Donc tout \((x,y)\) s’écrit ainsi, et la famille est génératrice.
\end{exemple}

\subsection{Sur-famille d’une famille génératrice}

\begin{proposition}
Toute famille contenant une famille génératrice est génératrice.
\end{proposition}

\begin{proof}
Soit \(\mathcal{F}\) une famille génératrice de \(E\), et \(\mathcal{G}\) une famille contenant \(\mathcal{F}\).  
Tout vecteur de \(E\) est combinaison linéaire de vecteurs de \(\mathcal{F}\), donc aussi de vecteurs de \(\mathcal{G}\). Ainsi \(\mathcal{G}\) est génératrice de \(E\).
\end{proof}

\begin{remarque}
Ajouter des vecteurs à une famille génératrice ne lui fait jamais perdre son caractère générateur. En revanche, cela peut introduire de la redondance.
\end{remarque}

\section{Familles libres}

\subsection{Définition}

\begin{definition}
Une famille finie \((u_1,\dots,u_p)\) de vecteurs de \(E\) est dite \textbf{libre} si la seule combinaison linéaire nulle de ces vecteurs est la combinaison triviale :
\[
\lambda_1u_1+\cdots+\lambda_pu_p=0_E
\Longrightarrow
\lambda_1=\cdots=\lambda_p=0.
\]
\end{definition}

\begin{definition}
Une famille qui n’est pas libre est dite \textbf{liée}.
\end{definition}

\begin{idee}
Une famille libre est une famille dans laquelle aucun vecteur n’est \textbf{de trop}.
\end{idee}

\begin{exemple}
Dans \(\R^2\), la famille
\[
\big((1,0),(0,1)\big)
\]
est libre.

En effet, si
\[
\lambda(1,0)+\mu(0,1)=(0,0),
\]
alors
\[
(\lambda,\mu)=(0,0),
\]
donc \(\lambda=0\) et \(\mu=0\).
\end{exemple}

\begin{exemple}
Dans \(\R^2\), la famille
\[
\big((1,0),(2,0)\big)
\]
est liée, car
\[
2(1,0)-(2,0)=(0,0)
\]
avec des coefficients non tous nuls.
\end{exemple}

\subsection{Cas d’une famille à un seul vecteur}

\begin{proposition}
La famille \((u)\) est libre si et seulement si \(u\neq 0_E\).
\end{proposition}

\begin{proof}
Si \(u=0_E\), alors
\[
1\cdot u=0_E
\]
avec un coefficient non nul, donc la famille est liée.

Réciproquement, si \(u\neq 0_E\) et si
\[
\lambda u=0_E,
\]
alors la propriété générale
\[
\lambda u=0_E \Longrightarrow \lambda=0 \text{ ou } u=0_E
\]
montre que \(\lambda=0\). La famille est donc libre.
\end{proof}

\subsection{Caractérisation d’une famille liée}

\begin{proposition}
Une famille finie \((u_1,\dots,u_p)\) est liée si et seulement s’il existe des scalaires \(\lambda_1,\dots,\lambda_p\), non tous nuls, tels que
\[
\lambda_1u_1+\cdots+\lambda_pu_p=0_E.
\]
\end{proposition}

\begin{proof}
C’est exactement la négation de la définition d’une famille libre.
\end{proof}

\subsection{Premiers critères pratiques}

\begin{proposition}
Si une famille contient le vecteur nul, alors elle est liée.
\end{proposition}

\begin{proof}
Supposons par exemple que \(u_1=0_E\). Alors
\[
1\cdot u_1+0\cdot u_2+\cdots+0\cdot u_p=0_E
\]
avec des coefficients non tous nuls. La famille est donc liée.
\end{proof}

\begin{proposition}
Si deux vecteurs d’une famille sont égaux, alors la famille est liée.
\end{proposition}

\begin{proof}
Si \(u_i=u_j\) avec \(i\neq j\), alors
\[
u_i-u_j=0_E,
\]
ce qui fournit une relation linéaire non triviale.
\end{proof}

\begin{proposition}
Si un vecteur d’une famille est combinaison linéaire des autres, alors la famille est liée.
\end{proposition}

\begin{proof}
Supposons par exemple
\[
u_p=\lambda_1u_1+\cdots+\lambda_{p-1}u_{p-1}.
\]
Alors
\[
\lambda_1u_1+\cdots+\lambda_{p-1}u_{p-1}-u_p=0_E,
\]
ce qui est une relation linéaire non triviale. La famille est donc liée.
\end{proof}

\subsection{Caractérisation opérationnelle d’une famille liée}

\begin{theoreme}
Soit \((u_1,\dots,u_p)\) une famille finie avec \(p\ge 2\).  
Cette famille est liée si et seulement si l’un des vecteurs est combinaison linéaire des autres.
\end{theoreme}

\begin{proof}
Si l’un des vecteurs est combinaison linéaire des autres, la proposition précédente montre que la famille est liée.

Réciproquement, supposons la famille liée. Il existe donc des scalaires \(\lambda_1,\dots,\lambda_p\), non tous nuls, tels que
\[
\lambda_1u_1+\cdots+\lambda_pu_p=0_E.
\]
L’un au moins des \(\lambda_i\) est non nul. Quitte à réordonner, supposons \(\lambda_p\neq 0\). Alors
\[
u_p=-\frac{\lambda_1}{\lambda_p}u_1-\cdots-\frac{\lambda_{p-1}}{\lambda_p}u_{p-1}.
\]
Donc \(u_p\) est combinaison linéaire des autres.
\end{proof}

\begin{remarque}
C’est souvent le critère le plus concret pour comprendre \emph{pourquoi} une famille est liée.
\end{remarque}

\section{Sous-familles d’une famille libre}

\begin{proposition}
Toute sous-famille d’une famille libre est libre.
\end{proposition}

\begin{proof}
Soit \((u_1,\dots,u_p)\) une famille libre, et extrayons-en une sous-famille, par exemple
\[
(u_{i_1},\dots,u_{i_r}).
\]
Supposons
\[
\lambda_1u_{i_1}+\cdots+\lambda_ru_{i_r}=0_E.
\]
On complète cette relation par des coefficients nuls sur les autres vecteurs de la famille initiale. On obtient une relation linéaire nulle sur \((u_1,\dots,u_p)\). Comme cette famille est libre, tous les coefficients sont nuls. En particulier,
\[
\lambda_1=\cdots=\lambda_r=0.
\]
La sous-famille est donc libre.
\end{proof}

\begin{remarque}
Retirer des vecteurs à une famille libre ne lui fait jamais perdre sa liberté.
\end{remarque}

\section{Caractérisation fondamentale des familles libres}

\begin{theoreme}
Une famille finie \((u_1,\dots,u_p)\) est libre si et seulement si chaque vecteur \(x\in \Vect(u_1,\dots,u_p)\) admet une unique écriture sous la forme
\[
x=\lambda_1u_1+\cdots+\lambda_pu_p.
\]
\end{theoreme}

\begin{proof}
Supposons d’abord la famille libre.

Soit \(x\in \Vect(u_1,\dots,u_p)\). Par définition, \(x\) admet au moins une écriture comme combinaison linéaire de la famille.

Montrons l’unicité. Supposons
\[
x=\lambda_1u_1+\cdots+\lambda_pu_p
\]
et
\[
x=\mu_1u_1+\cdots+\mu_pu_p.
\]
En soustrayant,
\[
(\lambda_1-\mu_1)u_1+\cdots+(\lambda_p-\mu_p)u_p=0_E.
\]
Comme la famille est libre,
\[
\lambda_1-\mu_1=\cdots=\lambda_p-\mu_p=0,
\]
donc
\[
\lambda_1=\mu_1,\dots,\lambda_p=\mu_p.
\]
L’écriture est donc unique.

Réciproquement, supposons que tout vecteur de \(\Vect(u_1,\dots,u_p)\) admette une unique écriture dans la famille. En particulier, \(0_E\in \Vect(u_1,\dots,u_p)\), et on a toujours
\[
0_E=0u_1+\cdots+0u_p.
\]
Si
\[
\lambda_1u_1+\cdots+\lambda_pu_p=0_E,
\]
alors \(0_E\) aurait deux écritures. Par unicité, on doit avoir
\[
\lambda_1=\cdots=\lambda_p=0.
\]
La famille est donc libre.
\end{proof}

\begin{remarque}
Cette propriété est essentielle : la liberté correspond exactement à l’unicité des coordonnées.
\end{remarque}

\section{Bases}

\subsection{Définition}

\begin{definition}
Une \textbf{base} d’un espace vectoriel \(E\) est une famille de vecteurs de \(E\) qui est à la fois :
\begin{itemize}
    \item \textbf{libre},
    \item \textbf{génératrice}.
\end{itemize}
\end{definition}

\begin{idee}
Une base est une famille qui contient \textbf{juste ce qu’il faut} :
\begin{itemize}
    \item assez de vecteurs pour engendrer tout l’espace ;
    \item pas trop de vecteurs pour éviter les redondances.
\end{itemize}
\end{idee}

\begin{exemple}
La famille
\[
\big((1,0),(0,1)\big)
\]
est une base de \(\R^2\). C’est la \textbf{base canonique} de \(\R^2\).
\end{exemple}

\begin{exemple}
La famille
\[
\big((1,1),(1,-1)\big)
\]
est une base de \(\R^2\) : elle est génératrice, et elle est libre.
\end{exemple}

\subsection{Une famille libre est une base de son espace engendré}

\begin{proposition}
Si \((u_1,\dots,u_p)\) est une famille libre, alors c’est une base de
\[
\Vect(u_1,\dots,u_p).
\]
\end{proposition}

\begin{proof}
Par définition, la famille \((u_1,\dots,u_p)\) engendre \(\Vect(u_1,\dots,u_p)\). Comme elle est libre, elle est donc à la fois libre et génératrice de cet espace : c’est une base.
\end{proof}

\begin{remarque}
C’est un résultat très utile : dès qu’on montre qu’une famille est libre, on obtient automatiquement une base du sous-espace qu’elle engendre.
\end{remarque}

\subsection{Base canonique}

\begin{definition}
Dans \(\K^n\), on appelle \textbf{base canonique} la famille
\[
(e_1,\dots,e_n),
\]
où
\[
e_1=(1,0,\dots,0),\quad e_2=(0,1,\dots,0),\quad \dots,\quad e_n=(0,\dots,0,1).
\]
\end{definition}

\begin{proposition}
La base canonique de \(\K^n\) est bien une base de \(\K^n\).
\end{proposition}

\begin{proof}
Soit \((x_1,\dots,x_n)\in \K^n\). On a
\[
(x_1,\dots,x_n)=x_1e_1+\cdots+x_ne_n.
\]
La famille est donc génératrice.

Supposons maintenant
\[
\lambda_1e_1+\cdots+\lambda_ne_n=0.
\]
Alors le membre de gauche est
\[
(\lambda_1,\dots,\lambda_n),
\]
donc
\[
(\lambda_1,\dots,\lambda_n)=(0,\dots,0),
\]
ce qui implique
\[
\lambda_1=\cdots=\lambda_n=0.
\]
La famille est libre. C’est donc une base.
\end{proof}

\begin{exemple}
Dans \(\R^3\), la base canonique est
\[
\big((1,0,0),(0,1,0),(0,0,1)\big).
\]
\end{exemple}

\section{Coordonnées dans une base}

\begin{theoreme}
Soit \(\mathcal{B}=(e_1,\dots,e_p)\) une base d’un espace vectoriel \(E\). Alors tout vecteur \(x\in E\) s’écrit de manière unique sous la forme
\[
x=\lambda_1e_1+\cdots+\lambda_pe_p.
\]
\end{theoreme}

\begin{proof}
Comme \(\mathcal{B}\) est génératrice, une telle écriture existe. Comme \(\mathcal{B}\) est libre, elle est unique.
\end{proof}

\begin{definition}
Les scalaires \(\lambda_1,\dots,\lambda_p\) sont appelés les \textbf{coordonnées} de \(x\) dans la base \(\mathcal{B}\).
\end{definition}

\begin{exemple}
Dans la base canonique de \(\R^3\),
\[
(2,-1,5)=2e_1-e_2+5e_3.
\]
Ses coordonnées sont donc
\[
(2,-1,5).
\]
\end{exemple}

\begin{exemple}
Dans la base \(\big((1,1),(1,-1)\big)\) de \(\R^2\), cherchons les coordonnées de \((3,1)\).

On résout
\[
(3,1)=\lambda(1,1)+\mu(1,-1),
\]
soit
\[
\begin{cases}
\lambda+\mu=3,\\
\lambda-\mu=1.
\end{cases}
\]
On obtient
\[
\lambda=2,\qquad \mu=1.
\]
Donc
\[
(3,1)=2(1,1)+1(1,-1).
\]
Les coordonnées de \((3,1)\) dans cette base sont \((2,1)\).
\end{exemple}

\section{Critères pratiques pour tester la liberté}

\subsection{Dans \(\K^2\)}

\begin{proposition}
Deux vecteurs \((u,v)\) d’un espace vectoriel sont libres si et seulement si l’un n’est pas multiple de l’autre.
\end{proposition}

\begin{proof}
Si \(v=\lambda u\), alors
\[
\lambda u-v=0,
\]
donc la famille est liée.

Réciproquement, si la famille est liée, il existe \((a,b)\neq (0,0)\) tel que
\[
au+bv=0.
\]
Si \(b\neq 0\), alors
\[
v=-\frac{a}{b}u.
\]
Sinon \(a\neq 0\) et
\[
u=0,
\]
qui est aussi un multiple de \(v\). Ainsi l’un des deux vecteurs est multiple de l’autre.
\end{proof}

\begin{exemple}
Dans \(\R^2\), les vecteurs
\[
(1,2)\quad \text{et} \quad (2,4)
\]
sont liés, car
\[
(2,4)=2(1,2).
\]
\end{exemple}

\begin{exemple}
Les vecteurs
\[
(1,2)\quad \text{et} \quad (2,3)
\]
sont libres, car aucun n’est multiple de l’autre.
\end{exemple}

\subsection{Méthode générale}

\begin{methode}
Pour tester si une famille \((u_1,\dots,u_p)\) est libre :
\begin{enumerate}
    \item on écrit
    \[
    \lambda_1u_1+\cdots+\lambda_pu_p=0_E ;
    \]
    \item on traduit cette égalité en système ;
    \item on montre que la seule solution est
    \[
    \lambda_1=\cdots=\lambda_p=0.
    \]
\end{enumerate}
\end{methode}

\begin{exemple}
Étudions la liberté de la famille
\[
\big((1,0,1),(0,1,1),(1,1,0)\big)
\]
dans \(\R^3\).

On suppose
\[
\lambda(1,0,1)+\mu(0,1,1)+\nu(1,1,0)=(0,0,0).
\]
Cela donne
\[
(\lambda+\nu,\mu+\nu,\lambda+\mu)=(0,0,0),
\]
soit
\[
\begin{cases}
\lambda+\nu=0,\\
\mu+\nu=0,\\
\lambda+\mu=0.
\end{cases}
\]
Des deux premières équations, on tire
\[
\lambda=-\nu,\qquad \mu=-\nu.
\]
En remplaçant dans la troisième :
\[
-\nu-\nu=0 \Longrightarrow \nu=0.
\]
Donc \(\lambda=\mu=\nu=0\). La famille est libre.
\end{exemple}

\section{Critères pratiques pour tester le caractère générateur}

\begin{methode}
Pour montrer qu’une famille \((u_1,\dots,u_p)\) est génératrice d’un espace \(E\), on prend un vecteur arbitraire \(x\in E\) et on cherche à résoudre
\[
x=\lambda_1u_1+\cdots+\lambda_pu_p.
\]
Si cette équation est toujours résoluble, alors la famille est génératrice.
\end{methode}

\begin{exemple}
Montrons que
\[
\big((1,0,1),(0,1,1),(1,1,0)\big)
\]
engendre \(\R^3\).

Soit \((x,y,z)\in\R^3\). Cherchons \(\lambda,\mu,\nu\in\R\) tels que
\[
\lambda(1,0,1)+\mu(0,1,1)+\nu(1,1,0)=(x,y,z).
\]
Cela donne
\[
\begin{cases}
\lambda+\nu=x,\\
\mu+\nu=y,\\
\lambda+\mu=z.
\end{cases}
\]
En résolvant :
\[
\lambda=\frac{x-y+z}{2},\qquad
\mu=\frac{-x+y+z}{2},\qquad
\nu=\frac{x+y-z}{2}.
\]
On peut donc toujours trouver une solution. La famille est génératrice de \(\R^3\).
\end{exemple}

\subsection{Lien avec les systèmes linéaires}

\begin{remarque}
Tester si une famille est libre revient à résoudre un \textbf{système homogène} :
\[
\lambda_1u_1+\cdots+\lambda_pu_p=0_E.
\]

Tester si une famille est génératrice revient à résoudre un \textbf{système avec second membre arbitraire} :
\[
x=\lambda_1u_1+\cdots+\lambda_pu_p.
\]

Ainsi, les notions de liberté et de génération sont étroitement liées à la résolution de systèmes linéaires.
\end{remarque}

\section{Extraction d’une base à partir d’une famille génératrice finie}

\begin{theoreme}
De toute famille génératrice finie d’un espace vectoriel, on peut extraire une base de cet espace.
\end{theoreme}

\begin{proof}
Soit \((u_1,\dots,u_p)\) une famille génératrice de \(E\).

Si cette famille est libre, c’est déjà une base, et il n’y a rien à faire.

Sinon, elle est liée. D’après la caractérisation opérationnelle des familles liées, l’un des vecteurs est combinaison linéaire des autres. En le retirant, on obtient une famille encore génératrice de \(E\), car le vecteur retiré pouvait déjà être fabriqué à l’aide des autres.

Si la nouvelle famille est libre, on a une base. Sinon, on recommence.

Comme la famille initiale est finie, ce processus s’arrête après un nombre fini d’étapes. On obtient alors une famille libre et génératrice : c’est une base de \(E\).
\end{proof}

\begin{remarque}
Ce résultat est fondamental. Il montre qu’une famille génératrice contient toujours une version plus efficace, sans redondance.
\end{remarque}

\begin{exemple}
Dans \(\R^3\), la famille
\[
\big((1,0,0),(0,1,0),(1,1,0),(0,0,1)\big)
\]
est génératrice. Le vecteur \((1,1,0)\) vaut
\[
(1,1,0)=(1,0,0)+(0,1,0),
\]
donc il est redondant. En le retirant, on obtient la base
\[
\big((1,0,0),(0,1,0),(0,0,1)\big).
\]
\end{exemple}

\section{Bases de quelques espaces classiques}

\subsection{Bases de \(\K^n\)}

\begin{exemple}
Une autre base de \(\R^2\) est
\[
\big((1,1),(1,-1)\big).
\]
Une autre base de \(\R^3\) est
\[
\big((1,0,0),(1,1,0),(1,1,1)\big).
\]
\end{exemple}

\subsection{Bases de polynômes}

\begin{exemple}
Dans l’espace des polynômes de degré inférieur ou égal à \(2\), une base naturelle est
\[
(1,X,X^2).
\]
En effet, tout polynôme \(P(X)\) de degré \(\le 2\) s’écrit
\[
P(X)=a+bX+cX^2
\]
de manière unique.
\end{exemple}

\begin{proposition}
La famille \((1,X,X^2)\) est libre dans \(\K[X]\).
\end{proposition}

\begin{proof}
Supposons
\[
\lambda\cdot 1+\mu X+\nu X^2=0
\]
comme polynôme nul. Cela signifie que pour tout \(X\),
\[
\lambda+\mu X+\nu X^2=0.
\]
Un polynôme nul a tous ses coefficients nuls, donc
\[
\lambda=\mu=\nu=0.
\]
La famille est libre.
\end{proof}

\subsection{Exemple dans un espace de fonctions}

\begin{exemple}
Dans l’espace des fonctions de \(\R\) dans \(\R\), la famille
\[
(1,\mathrm{id})
\]
où \(\mathrm{id}(x)=x\), est libre.

En effet, si
\[
\lambda\cdot 1+\mu\,\mathrm{id}=0
\]
comme fonction, alors pour tout \(x\in\R\),
\[
\lambda+\mu x=0.
\]
En prenant \(x=0\), on obtient \(\lambda=0\). Puis en prenant \(x=1\), on obtient \(\mu=0\). La famille est donc libre.
\end{exemple}

\begin{exemple}
Dans l’espace des fonctions polynomiales de degré \(\le 1\), la famille
\[
(1,\mathrm{id})
\]
est une base.
\end{exemple}

\subsection{Base d’un sous-espace}

\begin{exemple}
Dans \(\R^3\), considérons le plan
\[
F=\{(x,y,z)\in\R^3\mid x+y+z=0\}.
\]
Tout vecteur de \(F\) s’écrit
\[
(x,y,z)=(x,y,-x-y)=x(1,0,-1)+y(0,1,-1).
\]
Donc
\[
F=\Vect\big((1,0,-1),(0,1,-1)\big).
\]
Comme ces deux vecteurs ne sont pas colinéaires, ils sont libres. Ainsi
\[
\big((1,0,-1),(0,1,-1)\big)
\]
est une base de \(F\).
\end{exemple}

\section{Méthode d’analyse-synthèse dans ce chapitre}

\begin{remarque}
La méthode d’analyse-synthèse est très utile pour :
\begin{itemize}
    \item déterminer si une famille est génératrice ;
    \item trouver les coordonnées d’un vecteur dans une base ;
    \item exhiber une base d’un sous-espace.
\end{itemize}
\end{remarque}

\begin{exemple}
Trouvons une base du sous-espace
\[
F=\{(x,y,z)\in\R^3\mid x-y+z=0\}.
\]

\textbf{Analyse :} si \((x,y,z)\in F\), alors
\[
x-y+z=0 \Longrightarrow x=y-z.
\]
Donc
\[
(x,y,z)=(y-z,y,z)=y(1,1,0)+z(-1,0,1).
\]
Ainsi tout vecteur de \(F\) est combinaison linéaire de
\[
(1,1,0)\quad \text{et} \quad (-1,0,1).
\]

\textbf{Synthèse :} réciproquement, tout vecteur de la forme
\[
\lambda(1,1,0)+\mu(-1,0,1)=(\lambda-\mu,\lambda,\mu)
\]
vérifie
\[
(\lambda-\mu)-\lambda+\mu=0.
\]
Donc
\[
F=\Vect\big((1,1,0),(-1,0,1)\big).
\]
Comme ces deux vecteurs ne sont pas colinéaires, ils sont libres. C’est donc une base de \(F\).
\end{exemple}

\section{Méthode générale : libre, génératrice, base}

\begin{methode}
Pour montrer qu’une famille \((u_1,\dots,u_p)\) est \textbf{libre}, on :
\begin{enumerate}
    \item écrit
    \[
    \lambda_1u_1+\cdots+\lambda_pu_p=0_E ;
    \]
    \item résout ;
    \item vérifie que la seule solution est la solution triviale.
\end{enumerate}

Pour montrer qu’elle est \textbf{génératrice}, on :
\begin{enumerate}
    \item prend un vecteur arbitraire \(x\) de l’espace ;
    \item cherche à écrire
    \[
    x=\lambda_1u_1+\cdots+\lambda_pu_p ;
    \]
    \item montre que c’est toujours possible.
\end{enumerate}

Pour montrer qu’elle est une \textbf{base}, on montre les deux propriétés.
\end{methode}

\section{Erreurs classiques}

\begin{itemize}
    \item Confondre \textbf{libre} et \textbf{génératrice}.
    \item Penser qu’une famille génératrice est forcément libre.
    \item Penser qu’une famille libre est forcément génératrice.
    \item Oublier qu’une base doit satisfaire les deux propriétés.
    \item Croire que deux vecteurs non égaux sont automatiquement libres.
    \item Croire que trois vecteurs de \(\R^3\) forment automatiquement une base.
    \item Pour tester la liberté, ne pas partir de
    \[
    \lambda_1u_1+\cdots+\lambda_pu_p=0.
    \]
    \item Pour tester le caractère générateur, ne pas partir d’un vecteur arbitraire de l’espace.
    \item Confondre \(\Vect(u_1,\dots,u_p)\) avec une base de cet espace.
    \item Confondre coordonnées dans la base canonique et coordonnées dans une autre base.
\end{itemize}

\section{Résumé du chapitre}

\begin{itemize}
    \item Une famille est \textbf{génératrice} si elle engendre tout l’espace.
    \item Une famille est \textbf{libre} si la seule relation linéaire nulle est la relation triviale.
    \item Une famille \textbf{liée} admet une relation linéaire non triviale.
    \item Toute sous-famille d’une famille libre est libre.
    \item Toute sur-famille d’une famille génératrice est génératrice.
    \item Une famille liée est exactement une famille dans laquelle un vecteur est combinaison linéaire des autres.
    \item Une \textbf{base} est une famille à la fois libre et génératrice.
    \item Une famille libre est une base de son espace engendré.
    \item Dans une base, tout vecteur admet une écriture unique.
    \item Les coefficients de cette écriture sont les \textbf{coordonnées} du vecteur dans la base.
    \item De toute famille génératrice finie, on peut extraire une base.
\end{itemize}

\section*{Exercices d’application immédiate}

\begin{enumerate}[label=\textbf{Exercice \arabic*.}]
    \item Dans \(\R^2\), dire si les familles suivantes sont libres, génératrices, ou des bases :
    \[
    \big((1,0),(0,1)\big),\qquad
    \big((1,0),(2,0)\big),\qquad
    \big((1,1),(1,-1)\big).
    \]

    \item Montrer que la famille
    \[
    \big((1,0,0),(0,1,0),(0,0,1)\big)
    \]
    est une base de \(\R^3\).

    \item Étudier la liberté de la famille
    \[
    \big((1,2),(2,4)\big).
    \]

    \item Montrer que la famille
    \[
    \big((1,1),(1,-1)\big)
    \]
    est une base de \(\R^2\).

    \item Dans \(\R^3\), déterminer si la famille
    \[
    \big((1,0,1),(0,1,1),(1,1,0)\big)
    \]
    est une base.

    \item Trouver une base du sous-espace
    \[
    F=\{(x,y,z)\in\R^3\mid x+y-z=0\}.
    \]

    \item Dans l’espace des polynômes de degré inférieur ou égal à \(2\), montrer que
    \[
    (1,X,X^2)
    \]
    est une base.

    \item Dans la base \(\big((1,1),(1,-1)\big)\) de \(\R^2\), déterminer les coordonnées de \((5,1)\).

    \item Montrer que toute sous-famille d’une famille libre est libre.

    \item Montrer que toute famille contenant une famille génératrice est génératrice.

    \item Montrer qu’une famille libre est une base de son espace engendré.

    \item Dans l’espace des fonctions de \(\R\) dans \(\R\), montrer que la famille \((1,\mathrm{id})\) est libre.
\end{enumerate}

\end{document}